\documentclass{article}
\usepackage[T1]{fontenc}
\usepackage{lmodern}
\usepackage{graphicx}
\usepackage{amsmath,amssymb}
\usepackage[a4paper,margin=2cm]{geometry}
\usepackage[absolute,overlay]{textpos}
\setlength{\TPHorizModule}{1mm}
\setlength{\TPVertModule}{1mm}
\usepackage[indonesian]{babel}
\usepackage{hyperref}
\setlength{\emergencystretch}{2em}
% unit: Halaman 20

\begin{document}

% === Halaman 20 (python/native) ===

20

Contoh 2: Diberikan cipherteks berikut ini, yang dihasilkan dari cipher abjad-Tunggal. 

Lakukan dekripsi untuk mengungkap plainteksnya.

 LIVITCSWPIYVEWHEVSRIQMXLEYVEOIEWHRXEXIPFEMVEWHKVSTYLXZIX

LIKIIXPIJVSZEYPERRGERIMWQLMGLMXQERIWGPSRIHMXQEREKIETXMJT

PRGEVEKEITREWHEXXLEXXMZITWAWSQWXSWEXTVEPMRXRSJGSTVRIEYVI

EXCVMUIMWERGMIWXMJMGCSMWXSJOMIQXLIVIQIVIXQSVSTWHKPEGARCS

XRWIEVSWIIBXVIZMXFSJXLIKEGAEWHEPSWYSWIWIEVXLISXLIVXLIRGE

PIRQIVIIBGIIHMWYPFLEVHEWHYPSRRFQMXLEPPXLIECCIEVEWGISJKTV

WMRLIHYSPHXLIQIMYLXSJXLIMWRIGXQEROIVFVIZEVAEKPIEWHXEAMWY

EPPXLMWYRMWXSGSWRMHIVEXMSWMGSTPHLEVHPFKPEZINTCMXIVJSVLMR

SCMWMSWVIRCIGXMWYMX

\end{document}
